\chapter{数学推理}

\section{数学综合题}
\noindent\textbf{要点与方法.}\enspace
离散结构题可先寻找同余、奇偶、总量或关联次数等不变量，再把可行性分成必要条件与具体构造。双重计数连接局部重叠与整体规模；阶乘整除常可转为素因子的指数或素因子所在区间的问题。

几何优化先选坐标、写目标，并核对边界与小规模特例；在凸可行域上，严格凸性可用于证明最优点唯一。几何概率需保留变换后的密度和积分区域，等分布也不意味着固定间隔的事件独立。

分析题要标明定义域与正则条件。倒数变换、减去多项式可简化方程；差商序列用于否定可微性，有限前缀加统一尾界用于处理稠密跳点。离散分布的概率积分变换还需处理原子。
\par\smallskip



\problem{A14}{环面上的皇后排列}{同余不变量；构造}
把 $n\times n$ 棋盘的两对边分别接合，行列编号为模 $n$ 的余数。要求放置 $n$ 个皇后，每行每列各一个，且所占格 $(i,j)$ 的两种对角线编号 $i+j$ 与 $i-j$ 也各不重复。判定所有可行的正整数 $n$，并构造可行排列。
\solution 记所选列为置换 $\pi(i)$，$i=0,\ldots,n-1$。若 $i+\pi(i)$ 也为置换，则求和模 $n$ 得
\[
 \frac{n(n-1)}2\equiv n(n-1)\equiv0\pmod n,
\]
故 $n$ 必须为奇数。再令 $S=\sum_{i=0}^{n-1}i^2$。两种对角线都为置换，所以平方和满足
\[
 2S\equiv\sum_i\bigl((i+\pi(i))^2+(i-\pi(i))^2\bigr)
 =4S\pmod n.
\]
因此 $n\mid2S=n(n-1)(2n-1)/3$。若 $3\mid n$，则 $(n-1)(2n-1)$ 不被 $3$ 整除，此整除关系不可能成立。必要条件为 $\gcd(n,6)=1$。
反过来，若此条件成立，取 $\pi(i)=2i\pmod n$。列编号、正对角线编号和负对角线编号分别为 $2i,3i,-i$，其乘数均与 $n$ 互素，所以全是置换。$n=1$ 也由此构造满足要求。
\insight 一阶求和排除偶数，二阶求和排除三的倍数；两个方向的对角线约束需要不同的不变量。

\problem{A17}{七人的共享安排}{关联计数；差集构造}
七人各持有三种物品，任意两人恰好共有一种；每一种物品至少由一人、至多由三人持有。证明总共恰有七种物品，而且每种都由三人持有。构造满足全部要求的安排。
\solution 令 $n_j$ 为恰由 $j$ 人持有的种数，$j=1,2,3$。按人和物品两种方式计数所有持有关系，得到
\[
 n_1+2n_2+3n_3=21.
\]
每对人恰共享一种，按物品计数这 $\binom72=21$ 个共享关系，得到 $n_2+3n_3=21$。两式相减得 $n_1+n_2=0$，因此 $n_1=n_2=0,n_3=7$。
将人编号为 $0,\ldots,6$，对每个 $i$ 设置一种由
\[
 \{i,i+1,i+3\}\pmod7
\]
三人共同持有的物品。每个人在三个平移集合中出现，故各有三种。基本集合 $\{0,1,3\}$ 中六个有向差恰为模 $7$ 的六个非零余数，所以给定有序人员对，它恰出现在一个平移集合中；于是任意两人恰共有一种。
\insight 同时数“人持有物品”的次数与“人员对共同持有物品”的次数，可以强制每种物品的持有人数；构造再用差的唯一性验证。

\problem{A19}{图上的校验恢复}{异或方程；连通分量}
有限无向简单图的每条边带一个二元值。已知每个顶点所有关联边的异或和，以及部分边的值。令未知边构成的图为 $H$，保留所有 $v$ 个顶点，未知边数为 $q$，连通分量数为 $c$（包括孤立点）。给出数据相容的充要条件、相容时的解数以及唯一恢复条件，并说明如何构造所有解。
\solution 先从每个顶点的校验值中异或掉已知边，得到剩余要求 $b_x$。在 $H$ 的任一连通分量中，未知边都被端点计数两次，所以必要条件是该分量内 $b_x$ 的异或和为零。
为证充分性，每个非平凡连通分量取一棵生成树。任意指定全部非树边的值，再把它们从顶点要求中消去。树上从叶子开始：叶子唯一的剩余边必须等于该叶子的校验值，确定后删去叶子并更新邻点。最后根的要求由分量总异或为零自动满足。孤立点恰要求 $b_x=0$。
生成森林有 $v-c$ 条边，因此非树边共有 $q-v+c$ 条。其值可以独立任选，树边随后唯一确定，故解数为
\[
 \boxed{2^{q-v+c}}.
\]
相容时唯一恢复当且仅当 $q=v-c$，即未知边图是森林。若有环，把环上全部位翻转不会改变任何顶点校验，也直接说明解不唯一。
\insight 唯一性由未知部分的图结构决定，而不是由原图是否有环决定。已观测边可能恰好打断所有剩余环。

\problem{A20}{三账户的余额调整}{模不变量；最短操作}
三个账户的初始余额为整数 $a_1,a_2,a_3$。每步选择一个账户减去 $2$，另两个各加 $1$；中间余额允许为负。何时能使三个余额相等？若可行，求最少操作步数并给出实现方案。
\solution 总额不变，所以共同终值必须为 $t=(a_1+a_2+a_3)/3$。一次操作使任意两账户之差改变 $0$ 或 $\pm3$，所以初始余额必须两两模 $3$ 同余。此条件也使 $t$ 为整数。
假设共做 $N$ 步，第 $i$ 个账户被选作减款账户 $n_i$ 次，则 $\sum_i n_i=N$，其最终余额为
\[
 a_i+N-3n_i=t.
\]
因此 $n_i=(a_i+N-t)/3\ge0$，特别地 $N\ge t-\min_i a_i$。令 $m=\min_i a_i$，选择
\[
 N=t-m,\qquad n_i=\frac{a_i-m}{3}.
\]
在同余条件下这些是非负整数，且其和恰为 $N$。按任意顺序分别选择账户 $i$ 共 $n_i$ 次，即可达到三个 $t$。所以同余条件充分，而且最短步数就是 $t-m$。三余额本来相等时公式给出零步。
\insight 不变量给出可达障碍；把每类操作次数列为未知数，则能同时求出可达方案和最短长度。允许负余额是方案可以任意排序的必要规则。

\problem{A22}{中心二项式系数的二进赋值}{阶乘赋值；二进制位数}
对正整数 $n$，求 $\binom{2n}{n}$ 中因子 $2$ 的精确指数，并判定哪些 $n$ 使该数是偶数但不是 $4$ 的倍数。用二进制表示解释答案。
\solution 每个偶数给阶乘贡献至少一个 $2$，每个 $4$ 的倍数再贡献一个，依此类推，故
\[
 v_2(n!)=\sum_{j\ge1}\left\lfloor\frac n{2^j}\right\rfloor.
\]
写 $n=\sum_{k\ge0}b_k2^k$，$b_k\in\{0,1\}$，交换有限求和得到
\[
 v_2(n!)=\sum_k b_k(2^k-1)=n-s_2(n),
 \qquad s_2(n)=\sum_k b_k.
\]
因此
\[
 v_2\binom{2n}{n}
 =2n-s_2(2n)-2(n-s_2(n))=s_2(n),
\]
其中 $s_2(2n)=s_2(n)$，因为乘二只在末尾补零。正整数至少有一个非零二进制位，所以中心二项式系数恒为偶数；它不被 $4$ 整除当且仅当 $s_2(n)=1$，即 $n$ 是 $2$ 的非负整数次幂。
例如 $n=13=(1101)_2$ 时精确指数为 $3$，无需展开庞大的二项式系数。
\insight 阶乘的整除问题可以压缩成进制数位问题；精确赋值比仅证明偶性提供更多信息。

\problem{A15}{分界两侧的最短时间}{几何优化；折射关系}
平面以直线 $x=0$ 为界，左、右两侧行进速度分别为 $v_1,v_2>0$。从 $A=(-a,0)$ 到 $B=(b,h)$，其中 $a,b,h>0$。限定路径由两条直线段组成，并在 $C=(0,y)$ 处一次跨界。求最优交点的方程，证明其唯一，并给出等速情形与 $v_2\to\infty$ 时交点的极限。
\solution 总时间为
\[
 T(y)=\frac{\sqrt{a^2+y^2}}{v_1}
      +\frac{\sqrt{b^2+(h-y)^2}}{v_2},\qquad y\in\R.
\]
它在 $|y|\to\infty$ 时发散，而且
\[
 T''(y)=\frac{a^2}{v_1(a^2+y^2)^{3/2}}
 +\frac{b^2}{v_2(b^2+(h-y)^2)^{3/2}}>0.
\]
故存在唯一最小值。又 $T'(0)<0$、$T'(h)>0$，最优点 $y_*$ 位于 $(0,h)$，满足
\[
 \boxed{\frac{y_*}{v_1\sqrt{a^2+y_*^2}}
 =\frac{h-y_*}{v_2\sqrt{b^2+(h-y_*)^2}}}.
\]
若 $\theta_1,\theta_2$ 分别为两段与界面法线的夹角，这就是 $\sin\theta_1/v_1=\sin\theta_2/v_2$。
当 $v_1=v_2$ 时，方程两边均正，平方约分得 $by_*=a(h-y_*)$，所以 $y_*=ah/(a+b)$，路径恰是直线 $AB$。固定其他参数而 $v_2\to\infty$，方程右侧至多 $1/v_2\to0$，且 $0<y_*<h$，故必有 $y_*\to0$：应尽量减少慢速区域中的路程。
\insight 先用严格凸性证明驻点就是唯一全局最优，再求交点方程；无需为一般速度比强求复杂的根式。

\problem{A16}{折绳切点的优化}{展开坐标；最小最大值}
长为 $L>0$ 的绳子连续对折 $n\ge1$ 次，每次把当前右半段折向左半段。最后短段长 $h=L/2^n$，在距左端 $th$ 处一刀切断全部层，$0<t<1$。选择 $t$，使展开后最长一段尽可能短。分别处理 $n=1$ 与 $n\ge2$，并判断能否使所有段等长。
\solution 令 $N=2^n$。原绳上的 $N$ 个长为 $h$ 的小区间交替正向、反向覆盖短段，因此切点依次为
\[
 th,\ (2-t)h,\ (2+t)h,\ (4-t)h,\ldots,(N-t)h.
\]
两端段长为 $th$；内部相邻切点距离交替为 $2(1-t)h$ 和 $2th$。
当 $n=1$ 时，只有一个内部段，最长段为 $h\max\{t,2(1-t)\}$。两项相等时唯一达到最小值，即 $t=2/3$；三段均为 $2h/3=L/3$。
当 $n\ge2$ 时，两类内部段都出现，故最长段为
\[
 2h\max\{t,1-t\}\ge h.
\]
唯一最优是 $t=1/2$，最长段为 $h=L/2^n$。此时两端段只有 $h/2$，其余 $N-1$ 段均为 $h$。任何 $t>0$ 下端段 $th$ 与出现的内部段 $2th$ 都不等，因此不可能把全部段切成同样长。
\insight 不要只数切点；先把展开后的每类间距列全。小层数可能缺少某类内部间距，必须单独处理。

\problem{A18}{两随机点的街区距离}{几何概率；卷积}
在单位正方形内独立均匀选取两点 $(X_1,Y_1)$ 与 $(X_2,Y_2)$，记街区距离为 $D=|X_1-X_2|+|Y_1-Y_2|$。求 $0\le t\le1$ 时 $\Prob(D\le t)$，并求 $\Prob(D>1)$ 与 $\E D$。
\solution 对 $U=|X_1-X_2|$，单位正方形中两条对角三角形给出 $\Prob(U>u)=(1-u)^2$，所以密度为 $f_U(u)=2(1-u)$、$0<u<1$。$V=|Y_1-Y_2|$ 与 $U$ 独立同分布。
当 $0\le t\le1$ 时，积分区域是 $u,v\ge0,u+v\le t$，因此
\[
 \begin{aligned}
 \Prob(D\le t)
 &=4\int_0^t(1-u)\int_0^{t-u}(1-v)\,\dd v\,\dd u\\
 &=2t^2-\frac43t^3+\frac16t^4.
 \end{aligned}
\]
取 $t=1$ 得 $\Prob(D>1)=1-5/6=1/6$。另外
\[
 \E U=2\int_0^1u(1-u)\,\dd u=\frac13,
 \qquad\E D=\frac23.
\]
分布连续，故边界 $D=1$ 不影响严格与非严格尾概率。对于 $t>1$，积分区域会被 $u,v\le1$ 截断，不能继续使用上述多项式。
\insight 先求每个坐标差的分布，再用独立性卷积；距离均值则直接由线性性求出，无需先求完整分布。

\problem{A21}{相隔三项的首位数字}{等分布；联合频率}
从 $1,\ldots,N$ 均匀选取 $n$。当 $N\to\infty$ 时，求 $2^n$ 与 $2^{n+3}$ 的十进制首位数字同时为 $1$ 的概率极限，以及在 $2^n$ 首位为 $1$ 的条件下后者首位为 $1$ 的概率极限。两事件是否渐近独立？
\solution 令 $\alpha=\log_{10}2$。若其为有理数，将有某个 $2^q=10^p$，违背素因子分解，故 $\alpha$ 无理。序列 $\{n\alpha\}$ 在 $[0,1)$ 等分布：对任意非零整数 $h$，相应指数和是有界几何和，除以 $N$ 后趋于零，Weyl 判据即给出结论。
首位为 $1$ 等价于小数部分落在 $I=[0,\alpha)$。相隔三项对应圆周平移 $3\alpha=\log_{10}8$。对 $x\in I$，若 $x+3\alpha<1$，它不可能再落入 $I$；若越过 $1$，则要求
\[
 1-3\alpha\le x<\alpha.
\]
平移后的上界 $x<1-2\alpha$ 在 $x<\alpha$ 时自动成立，因为 $3\alpha<1$。所以联合极限为
\[
 4\alpha-1=\log_{10}(8/5),
\]
条件极限为 $(4\alpha-1)/\alpha$，约为 $0.6781$。边际极限为 $\alpha\approx0.3010$，联合极限约 $0.2041$ 不等于 $\alpha^2\approx0.0906$，故不渐近独立。
\insight 等分布给出单个时刻的均匀相位，却不使固定间隔的两个相位独立；联合频率要计算圆周区间的交集。

\problem{A23}{微分不等式与爆破时间}{倒数变换；尖锐界}
设 $a>0$，$f:[0,T)\to(0,\infty)$ 可微（零点取右导数），满足 $f(0)=a$ 和 $f'(x)\ge f(x)^2$。证明必有 $T\le1/a$，并给出在最长可能区间上满足条件的函数。进一步给出 $x<\min\{T,1/a\}$ 时 $f(x)$ 的下界。
\solution 令 $g=1/f$。正性保证该变换有定义，且
\[
 g'(x)=-\frac{f'(x)}{f(x)^2}\le-1.
\]
由中值定理，$g(x)-g(0)\le-x$，于是
\[
 0<\frac1{f(x)}\le\frac1a-x.
\]
若 $T>1/a$，可取 $x=1/a$ 得到矛盾。因此 $T\le1/a$。在右端为正的区间上取倒数，得到
\[
 f(x)\ge\frac{a}{1-ax}.
\]
函数 $f(x)=a/(1-ax)$ 在 $[0,1/a)$ 上为正且满足 $f'=f^2$，所以存在区间长度的界是尖锐的；它在右端发散，无法以有限正值可微地延拓过去。
注意一般满足不等式的函数可能更早爆破，或本来只在更短区间上给出；题设并不保证恰在 $1/a$ 处爆破，也不强制函数具有上述等号形式。
\insight 对二次增长的导数下界，倒数把非线性增长变成线性下降，从而直接控制最大存在时间。

\problem{A24}{二次扰动的加法方程}{Cauchy 方程；局部有界性}
求所有函数 $f:\R\to\R$，满足
\[
 f(x+y)=f(x)+f(y)+2xy\quad(x,y\in\R),
 \qquad f(1)=0,
\]
且 $f$ 在 $[0,1]$ 上有界。说明有界性在证明中的用途。
\solution 令 $g(x)=f(x)-x^2$，原式化为 $g(x+y)=g(x)+g(y)$，且 $g(1)=-1$。先得 $g(0)=0$、$g(-x)=-g(x)$，并由整数倍关系得到 $g(q)=-q$ 对一切有理数成立。
由于 $f$ 在 $[0,1]$ 有界，$g$ 亦有界；奇性进一步给出 $|g(x)|\le M$ 对 $|x|\le1$ 成立，可取 $M>0$。给定 $\varepsilon>0$，选整数 $k>M/\varepsilon$。若 $|h|<1/k$，则
\[
 |g(h)|=\frac{|g(kh)|}{k}\le\frac Mk<\varepsilon.
\]
故 $g$ 在零点连续，由加法性在每一点连续。以有理数 $q_j\to x$ 逼近，得 $g(x)=\lim_jg(q_j)=-x$。唯一答案为
\[
 \boxed{f(x)=x^2-x}.
\]
代回可直接核实方程、归一化与有界性。有界性用于把有理数上的结论延拓到实数；若去掉任何正则条件，加法方程存在非连续解，不能仅凭有理数计算就断言 $g(x)=-x$。
\insight 先减去能产生交叉项的二次函数，再处理加法方程；局部正则性只在从有理数推广到实数时使用。

\problem{A26}{振荡函数的局部正则性}{差商；振幅与频率}
设 $\alpha,\beta>0$，定义 $f(0)=0$，并对 $x\ne0$ 定义
\[
 f(x)=|x|^\alpha\sin(|x|^{-\beta}).
\]
判定 $f$ 在零点是否连续，何时在零点可微，以及何时导函数在零点连续。必要条件须用具体序列证明。
\solution 由 $|f(x)|\le|x|^\alpha$，对所有给定参数都在零点连续。差商绝对值至多 $|x|^{\alpha-1}$，所以 $\alpha>1$ 时 $f'(0)=0$。
反之，取正数序列使 $x_k^{-\beta}=2\pi k+\pi/2$ 与 $y_k^{-\beta}=2\pi k+3\pi/2$。相应差商为 $x_k^{\alpha-1}$ 与 $-y_k^{\alpha-1}$：$\alpha=1$ 时极限不同，$\alpha<1$ 时分别趋于正负无穷。因此可微当且仅当 $\alpha>1$。
在 $x>0$ 时，
\[
 f'(x)=\alpha x^{\alpha-1}\sin(x^{-\beta})
 -\beta x^{\alpha-\beta-1}\cos(x^{-\beta}).
\]
负半轴由偶函数的导数为奇函数得到对应式。若 $\alpha>\beta+1$，两项均趋于零，导函数连续。若 $1<\alpha\le\beta+1$，取 $x_k^{-\beta}=2\pi k$，得 $f'(x_k)=-\beta x_k^{\alpha-\beta-1}$，在等号时恒为 $-\beta$，在严格小于时发散，均不趋于 $f'(0)=0$。故导函数连续当且仅当 $\alpha>\beta+1$。
\insight 可微性由振幅除以距离控制，导函数连续性还受到振荡频率的影响，二者阈值不同。

\problem{A27}{单侧导数与凸性}{割线比较；可数异常集}
设 $f:\R\to\R$ 连续，每点都有有限右导数 $g(x)=f'_+(x)$，且 $g$ 单调不减。证明 $f$ 是凸函数，左导数处处存在且有限，并证明 $f$ 不可微的点至多可数。不得假设 $g$ 连续。
\solution 先证比较引理：连续函数 $h$ 若在 $[a,b)$ 上右导数非负，则 $h(b)\ge h(a)$。否则可选足够小的 $\varepsilon>0$，使 $q(x)=h(x)+\varepsilon x$ 仍有 $q(b)<q(a)$。$q$ 在紧区间的最大值必在某个 $t<b$ 取得；此时右差商非正，与 $q'_+(t)>0$ 矛盾。
将引理用于 $f(x)-g(a)x$ 和 $g(b)x-f(x)$，得
\[
 g(a)\le\frac{f(b)-f(a)}{b-a}\le g(b),\qquad a<b.
\]
因此对 $a<b<c$，前一段割线斜率不大于后一段割线斜率。整理这个不等式即为凸性不等式。
固定 $x$。单调函数 $g$ 的左极限 $g(x-)$ 存在，并被邻近两点的有限函数值夹住。对 $a<x$，比较引理还可用上界 $g(x-)$，因为 $g(t)\le g(x-)$ 对 $t<x$ 成立。于是
\[
 g(a)\le\frac{f(x)-f(a)}{x-a}\le g(x-).
\]
令 $a\uparrow x$，得 $f'_-(x)=g(x-)$。故不可微恰发生在 $g(x-)<g(x)$ 的点。给每个这样的点选一个有理数落在开区间 $(g(x-),g(x))$ 内；不同点的区间互不相交，所以这个选择给出到有理数集的单射。不可微点因此至多可数。
\insight 单调的单侧导数通过割线比较产生凸性；不可微点对应导数的正跳跃，而不需要导函数连续。

\problem{A28}{离散分布的随机化变换}{原子分布；逆分布耦合}
将有理数无重复枚举为 $r_1,r_2,\ldots$，设 $\Prob(X=r_j)=w_j>0$、$\sum_jw_j=1$，分布函数为 $F$。令 $V\sim\operatorname{Unif}(0,1)$ 与 $X$ 独立，定义
\[
 U=F(X-)+V\bigl(F(X)-F(X-)\bigr).
\]
证明 $U$ 在 $(0,1)$ 上均匀，并证明用广义逆 $Q(u)=\inf\{x:F(x)\ge u\}$ 可几乎处处恢复 $Q(U)=X$。为何不能直接使用 $F(X)$ 得到均匀变量？
\solution 在 $X=r_j$ 条件下，$U$ 在区间
\[
 I_j=(F(r_j-),F(r_j))
\]
上均匀，区间长度恰为 $w_j$。这些区间互不相交，长度之和为一，因此它们覆盖 $(0,1)$ 中除一个零测集外的全部点。对任意 Borel 集 $B\subset(0,1)$，
\[
 \Prob(U\in B)=\sum_jw_j\frac{|B\cap I_j|}{w_j}
 =|B|,
\]
这里 $|\cdot|$ 表示 Lebesgue 长度，故 $U$ 均匀。
若 $u\in I_j$，则任意 $x<r_j$ 都有 $F(x)\le F(r_j-)<u$，而 $F(r_j)>u$。所以 $Q(u)=r_j$，并由条件概率一的区间事件得到 $Q(U)=X$ 几乎处处。
$F(X)$ 只能落在可数集合 $\{F(r_j):j\ge1\}$ 中，故不可能服从连续均匀分布。辅助变量把每个原子的概率均匀铺开到其跳跃区间，正好消除了这个障碍。
\insight 连续分布的概率积分变换需要在原子处补充随机化；逆分布变换则能把同一均匀变量重新还原为原来的原子。

\par\Needspace{8\baselineskip}
\subsection*{小节练习}
\exercise{OA17}{四人的明信片}
四个人各有 $10$ 种不同的明信片。任意两人的收藏中，重复的种类不超过 $6$ 种。四人合起来至少有多少种不同的明信片？证明下界并给出达到下界的收藏方案。

\exercise{OA22}{阶乘整除的不等式}
设 $x,y$ 为正整数，且
\[
 \frac{x!\,y!}{x!+y!}
\]
为整数。证明 $3x-2y\ge2$。

\exercise{OA28}{稠密跳跃函数}
将所有有理数不重不漏地排列为 $r_0,r_1,r_2,\ldots$。定义
\[
 f(x)=\sum_{\{n\ge0:r_n<x\}}2^{-n},\qquad x\in\R.
\]
研究 $f$ 的连续性，指出全部连续点与不连续点，并求有理点处的左右极限。


\section{博弈问题}
\noindent\textbf{要点与方法.}\enspace
博弈题先列行动顺序、可见信息、支付与平局规则，再写单方偏离后的真实收益。私人信息或成交筛选会改变条件分布；优化报价前须先确定面对的是哪一类对象。个人均衡与总成本最优应分别计算。

零和混合策略可通过双方的收益保证证明最优；支持内等收益只产生候选，还要检查所有支持外行动。有限取物游戏可用强归纳分类局面，或设计每次都合法的对称回应，不能仅凭几个小例子断言周期。

构造题也要检验实施限制：匹配需核对邻集条件，同余组需检查相容性，称量编码除了区分状态，还必须满足每次两盘的平衡要求。
\par\smallskip



\problem{P01}{质量检测的价值}{逆向选择；信息设计}
卖方知道物品价值 $X\sim\operatorname{Unif}[0,1]$，买方估值为 $3X/2$。买方报一次价格 $p\ge0$，卖方当且仅当 $X\le p$ 时成交。买方可先支付费用 $c\ge0$ 做一次检测：他事先选择阈值 $t\in(0,1)$，随后仅得知 $X<t$ 或 $X\ge t$，再分别报价。求最优检测阈值、两种结果下的报价，并判断何时应购买检测。
\solution 无检测时，$0\le p\le1$ 的期望净交易收益为
\[
 \int_0^p(3x/2-p)\,\dd x=-p^2/4\le0;
\]
$p>1$ 时为 $3/4-p<0$。所以不交易可得最优收益零。
低质量结果下 $X$ 均匀于 $[0,t)$。报价 $p\le t$ 的收益非正；报 $p\ge t$ 时条件均值收益为 $3t/4-p<0$。故应不交易。
高质量结果下，只需考虑 $p\in[t,1]$，因为低于 $t$ 不成交、高于 $1$ 只会增加付款。将条件收益乘以该结果概率 $1-t$，得到无条件贡献
\[
 H_t(p)=\int_t^p(3x/2-p)\,\dd x
 =-\frac{p^2}{4}+tp-\frac{3t^2}{4}.
\]
这是关于 $p$ 的凹二次式，最优价格为 $p_t=\min\{2t,1\}$。检测带来的交易收益为
\[
 W(t)=\begin{cases}
 t^2/4,&0<t\le1/2,\\
 t-1/4-3t^2/4,&1/2\le t<1.
 \end{cases}
\]
第一段递增，第二段在 $t=2/3$ 处唯一最大。因此最优阈值为 $2/3$；低结果不买，高结果报价 $1$，毛收益为 $1/12$。扣除费用后，相对于不检测的改善为 $1/12-c$：费用低于 $1/12$ 时检测，高于时不检测，等于时两者无差异。
\insight 报价会筛选成交对象，检测又改变报价前的信息。先在每个信息分支优化报价，再优化检测规则，不能用无条件均值代替分支价值。


\problem{A07}{一次搜查两个位置}{零和博弈；覆盖策略}
甲把目标藏在三个位置之一，乙同时选择两个不同位置搜查。若查到位置 $i$，甲支付乙 $w_i$，其中 $(w_1,w_2,w_3)=(2,3,4)$；没查到则支付零。双方可随机化。求全部最优策略及博弈值。
\solution 甲以概率 $p=(6,4,3)/13$ 藏目标，则每个位置对期望支付的贡献均为 $w_ip_i=12/13$。任何两点搜查的收益都是 $24/13$，所以甲能将支付控制在此值。
乙以概率 $(1,5,7)/13$ 分别放弃位置 $1,2,3$，即选择其余两点。三个位置被搜查的概率为 $(12,8,6)/13$，乘以各自支付后均为 $24/13$。所以不论甲藏哪里，乙都保证获得这个值。两种保证相合，博弈值为
\[
 \boxed{V=24/13}.
\]
还需检查唯一性。上述乙的三个纯搜查动作都有正概率，因此若甲的其他策略也最优，三个两点搜查收益必须全等；否则加权平均为 $V$ 时至少一项超过 $V$。三个两项和全等迫使 $w_ip_i$ 全等，归一化后恰为上述 $p$。
类似地，甲的三个藏点都有正概率，所以乙的任何最优策略必须对各藏点给出相同支付 $V$，即搜查概率为 $V/w_i$，从而放弃各点的概率也唯一。故双方最优策略均唯一。
\insight 搜查动作是位置集合；应等化各集合的收益，并用另一方的覆盖概率提供独立的最优性证明。

\problem{A08}{新增抽查手段的影响}{混合策略；支持集变化}
检查者选择全面检查 $I$ 或不检查 $N$，另一方同时选择违规 $C$ 或守规 $H$。检查者的收益矩阵（行 $I,N$，列 $C,H$）为
\[
 \begin{pmatrix}1&-1\\-2&0\end{pmatrix},
\]
另一方的收益为其相反数。求混合均衡。随后允许检查者增加动作 $S$，它对两列的期望收益均为 $-1/4$；重新求全部最优策略和博弈值。
\solution 设违规概率为 $p$，两种检查动作的收益为 $2p-1$ 与 $-2p$。令其相等得 $p=1/4$，值为 $-1/2$。设全面检查概率为 $q$，对违规与守规的收益分别为 $3q-2$ 与 $-q$；令其相等得 $q=1/2$，同值为 $-1/2$。两方各纯策略收益已全部核验，故此为唯一混合均衡。
加入 $S$ 后，检查者纯用 $S$ 即保证 $-1/4$。另一方只要满足
\[
 2p-1\le-1/4,\qquad -2p\le-1/4,
\]
即 $p\in[1/8,3/8]$，就使所有检查动作收益不超过 $-1/4$。因此新博弈值是 $-1/4$，这些 $p$ 恰为全部最优违规概率。
检查者的唯一最优策略是纯用 $S$：面对 $p=1/4$，两个旧动作均仅得 $-1/2$，任何正比例混入旧动作都会使收益严格低于 $-1/4$。新增动作虽未改变另外两个动作之间的交点，却改变了最优支持集。
\insight 等化方程只在所选支持集内有效。加入动作后必须重新检查整张收益表，不能保留原混合比例不变。

\problem{A09}{固定减法的取石游戏}{递推分类；周期策略}
一堆有 $n\ge0$ 颗石子。两人轮流取走 $1$、$3$ 或 $4$ 颗，不得超过剩余数量；取完最后一颗者胜，轮到无石可取者负。判定全部必败局面，并给出从任一必胜局面出发的合法必胜首步。
\solution 记必败为 P、必胜为 N。初始七个状态 $n=0,\ldots,6$ 依次为
\[
 \mathrm{P,N,P,N,N,N,N}.
\]
例如 $n=2$ 只能留下必胜的 $1$，所以必败；$n=3,4$ 可以直接取完，$5,6$ 可分别取 $3,4$ 留下 $2$。
断言全部必败态满足 $n\bmod7\in\{0,2\}$。从余数 $0$ 取 $1,3,4$ 后落在 $6,4,3$；从余数 $2$ 则落在 $1,6,5$，均不在必败余数中。反过来，余数 $1,3,4,5,6$ 可分别取
\[
 1,\ 1,\ 4,\ 3,\ 4
\]
颗，落到余数 $0$ 或 $2$。在最小代表 $1,3,4,5,6$ 处这些取数也合法，因此对所有 $n$ 都合法。以较小局面作强归纳，两个方向共同证明分类及策略。
\insight 观察有限表只能提出周期猜想；完整证明必须检查必败态没有通向必败态的边，且每个其余状态都有这样的合法边。

\problem{A10}{异速服务器的拥挤均衡}{纯策略均衡；社会成本}
$n\ge2$ 名用户同时选择慢服务器或快服务器。若慢服务器有 $k$ 人，每人的成本为 $k$；快服务器有 $n-k$ 人，每人的成本为 $(n-k)/2$。用户只最小化自身成本，平手无差异。求全部纯策略纳什均衡的人数配置。对 $n=8$，再比较均衡与最小总成本配置。
\solution 当两台服务器都有人时，慢服务器用户不愿转向快服务器，当且仅当
\[
 k\le\frac{n-k+1}{2};
\]
快服务器用户不愿转向慢服务器，当且仅当 $(n-k)/2\le k+1$。所以均衡条件为
\[
 \boxed{n-2\le3k\le n+1}.
\]
$k=n$ 时有人转向空的快服务器可降低成本，故不均衡。$k=0$ 时只需检查快服务器用户：不转移当且仅当 $n/2\le1$，即 $n=2$；这也恰符合显示的条件。故在整数范围 $0\le k\le n$ 内，上式给出全部均衡；任意具有这种人数的用户划分都是均衡。
$n=8$ 时 $k=2$ 或 $3$。总成本为
\[
 C(k)=k^2+\frac{(8-k)^2}{2}
 =\frac32\left(k-\frac83\right)^2+\frac{64}{3}.
\]
整数最优为 $k=3$，成本 $43/2$；另一均衡 $k=2$ 的成本为 $22$。因此个人无改进机会不保证总成本最小，最差均衡与最优成本之比为 $44/43$。
\insight 单方转移时必须把目的地人数增加一。均衡条件与总成本优化是两个不同问题，应分别建立目标。

\problem{A11}{非均匀顾客的选址均衡}{同时博弈；中位数}
两家服务商同时在 $[0,1]$ 上选择位置，可以重合。顾客位置 $T$ 的密度为 $2t$；顾客选择较近的服务商，距离相同时两家平分其需求。每家最大化自己的需求份额，不收取价格。求全部纯策略纳什均衡，并计算均衡时顾客到服务商的平均距离。
\solution 顾客分布函数为 $F(t)=t^2$。若两家位置不同，记 $a<b$，左方份额为 $F((a+b)/2)$。左方把位置移到任意 $a'\in(a,b)$，分界点严格右移；由于密度在每个非空区间上的积分为正，份额严格增加。因此任何不同位置都不可能是均衡。
若两家重合于 $x$，各获 $1/2$。单方移到 $z<x$ 时获 $F((z+x)/2)$；移到 $z>x$ 时获 $1-F((z+x)/2)$。若 $F(x)>1/2$，从左侧充分接近 $x$ 可得到超过一半的份额；若 $F(x)<1/2$，从右侧接近可改进。因此唯一候选是中位数
\[
 m=1/\sqrt2.
\]
在该点，任意左偏离的份额严格小于 $F(m)=1/2$，任意右偏离严格小于 $1-F(m)=1/2$，不动仍得一半，所以双方都位于 $m$ 是唯一纯均衡。
均衡时路程为 $|T-m|$，因此
\[
 \begin{aligned}
 \E|T-m|&=\int_0^m2t(m-t)\,\dd t
       +\int_m^12t(t-m)\,\dd t\\
 &=\frac23-m+\frac23m^3=\frac{2-\sqrt2}{3}.
 \end{aligned}
\]
\insight 先排除分离位置，再检查重合位置的单方偏离。需求中位数决定竞争位置，不能用平均顾客位置代替。

\problem{A12}{保留价与密封拍卖}{私人价值；卖方决策}
两名风险中性买家的私人价值独立且均匀分布于 $[0,1]$。卖方先设保留价 $r\in[0,1]$，买家随后同时提交密封报价，也可不参与。最高的合格报价（至少为 $r$）获胜并支付自己的报价，其他人不支付；同价随机打破。构造对称均衡报价，计算卖方期望收入，并求最优 $r$。
\solution 对 $0<r<1$，考虑价值 $v<r$ 不参与、$v\ge r$ 报
\[
 b(v)=\frac{v^2+r^2}{2v}.
\]
$b(r)=r$，且在 $v>r$ 上严格递增。若价值 $v$ 的买家模拟类型 $t\in[r,1]$ 的报价，他获胜当且仅当对手价值小于 $t$（零概率同价不影响结果），期望收益为
\[
 (v-b(t))t=vt-\frac{t^2+r^2}{2}.
\]
该式关于 $t$ 严格凹，导数为 $v-t$。当 $v\ge r$ 时最优 $t=v$，收益 $(v^2-r^2)/2\ge0$；当 $v<r$ 时区间内最大值在 $t=r$，收益 $r(v-r)<0$，故不参与最优。低于保留价的报价等价于不参与，高于 $b(1)$ 的报价不会提高获胜概率却增加支付；所有偏离均已覆盖。
最大价值的密度为 $2v$，故卖方收入为
\[
 R(r)=\int_r^12v\,b(v)\,\dd v
 =\frac13+r^2-\frac43r^3.
\]
端点 $r=0$ 时 $b(v)=v/2$（定义 $b(0)=0$），$r=1$ 时几乎必然流拍，上式仍给出正确收入。由 $R'(r)=2r(1-2r)$，收入在 $r=1/2$ 唯一最大，为 $5/12$。
\insight 用“模拟另一类型”来检查报价偏离，可以同时推导收益与验证均衡；卖方收入还必须计入无人达到保留价的事件。

\problem{A13}{成排筹码的镜像策略}{对称配对；不变局面}
一排 $n\ge1$ 个位置各放一枚筹码。双方轮流拿走一枚，或拿走原始编号相邻且当前仍存在的两枚。空缺不会使原来不相邻的位置变成相邻；拿走最后筹码者胜。对任意 $n$，给出先手必胜策略并证明每次回应都合法。
\solution 若 $n$ 为奇数，先手拿走正中一枚；若 $n$ 为偶数，先手拿走正中的相邻两枚。余下位置关于原来中点反射对称，且中间有空缺。
以后对方每次拿走一枚或相邻两枚，先手就拿走它们的镜像。对方的合法动作不能跨越中间空缺，因此与镜像动作不重叠。对方行动前局面反射对称，所以镜像位置仍有筹码；反射也保持原始相邻关系，故回应合法。回应后恢复对称局面。
每轮双方的行动都减少筹码，过程有限。对方只要有合法动作，先手就有合法回应，因此不可能由对方拿走最后筹码。先手获胜。$n=1,2$ 时开局动作直接取完，也包括在策略内。
\insight 配对策略不仅要说明“对称”，还要保证对方动作与镜像动作不重合；先制造中心空缺正是为了解决这个问题。

\problem{CM01}{正则二分图的轮次安排}{Hall 条件；完美匹配分解}
有 $n$ 名员工和 $n$ 项任务。每名员工恰能执行 $d$ 项任务，每项任务也恰有 $d$ 名员工能执行，其中 $1\le d\le n$。证明能安排 $d$ 轮，每轮每名员工执行一项不同任务，而且每条允许的员工--任务配对恰使用一次。给出所用匹配定理的证明要点。
\solution 将允许配对画成二分图。对任意员工集合 $S$，其发出的边数为 $d|S|$，这些边都到达邻集 $N(S)$；每个任务最多接收 $d$ 条，故 $d|S|\le d|N(S)|$，即 $|N(S)|\ge|S|$。
Hall 定理说，此条件保证存在覆盖全部员工的匹配。可按左侧顶点数归纳证明：若有非空真子集 $S$ 满足等号，则分别匹配 $S$ 与 $N(S)$；在剩余图中，对任何员工集 $T$，原条件应用于 $S\cup T$ 给出 $|N(T)\setminus N(S)|\ge|T|$，所以剩余部分也可归纳匹配。
若每个非空真子集都满足严格不等式，任取一条边并删去其两个端点。剩余任意员工集的邻居至多减少一个，仍满足 Hall 条件，故可归纳匹配，再补回所选边。单个顶点是直接成立的基例。
因此原图存在一个完美匹配，作为第一轮。删掉其边后每个顶点的度数都变为 $d-1$；对剩余正则图重复，直到度数为零，得到 $d$ 个互不重叠的完美匹配。每轮覆盖所有人和任务，总共恰使用全部 $nd$ 条允许边一次。
\insight 局部度数条件通过边计数推出 Hall 条件；正则性又在删除一轮之后保留，从而把一次匹配提升为完整分解。

\problem{CM02}{非互素模数的同余构造}{广义中国剩余定理；扩展欧几里得}
给定正整数 $m,n$ 及整数 $a,b$，求同余组 $x\equiv a\pmod m$、$x\equiv b\pmod n$ 有解的充要条件，并给出可执行构造及解的周期。再求
\[
 x\equiv2\pmod6,\qquad x\equiv5\pmod9,
 \qquad x\equiv4\pmod5
\]
的全部整数解。
\solution 令 $g=\gcd(m,n)$。若有解，两式相减说明 $g\mid b-a$，这是必要条件。满足时写 $x=a+mt$，第二式等价于
\[
 \frac m g t\equiv\frac{b-a}{g}\pmod{n/g}.
\]
约分后的两个模系数互素，所以扩展欧几里得算法可求 $m/g$ 在模 $n/g$ 下的逆元，由此解出 $t$。若 $n/g=1$，则没有剩余限制，可取 $t=0$。因而条件也充分。
任意两个解的差同时被 $m,n$ 整除，故被 $\operatorname{lcm}(m,n)=mn/g$ 整除；反过来加上其任意整数倍仍是解。因此解恰为一个模 $mn/g$ 的剩余类。
例子中先写 $x=2+6t$，得到 $6t\equiv3\pmod9$，即 $2t\equiv1\pmod3$，所以 $t\equiv2\pmod3$、$x\equiv14\pmod{18}$。再写 $x=14+18k$，第三式给出 $3k\equiv0\pmod5$，即 $k\equiv0\pmod5$。全部解为
\[
 \boxed{x=14+90j,\qquad j\in\mathbb Z}.
\]
\insight 非互素模数首先需要检查相容性；约去公因子以后才能求逆，不能直接把不可逆的系数除掉。

\problem{CM03}{十二球的预定称量}{信息下界；三元编码}
十二枚外观相同的球中恰有一枚质量异常，可能较重或较轻，正常球质量相同。用无砝码天平，每次两盘放相同枚数的球，设计三次称量，且三次安排须事先固定、不随结果改变，以确定异常球及轻重。证明两次不够，并证明同类三次预定方案不能处理十三枚球。
\solution 每次结果为左重、平衡、右重，分别记 $1,0,-1$。十二球轻重共 $24$ 种状态，而两次只有 $3^2=9$ 种结果，所以至少三次。
以下矩阵每列对应一枚球，每行对应一次称量；$1$ 放左盘、$-1$ 放右盘、$0$ 不称：
\[
 \begin{array}{c|rrrrrrrrrrrr}
 &1&2&3&4&5&6&7&8&9&10&11&12\\\hline
 1&0&0&0&0&-1&-1&-1&1&1&1&-1&1\\
 2&0&-1&-1&-1&1&1&1&0&0&0&-1&1\\
 3&-1&1&0&-1&1&0&-1&-1&0&1&1&0
 \end{array}
\]
每行正负项各四个，所以称量合法。每列非零，且任意两列既不相同也不互为相反数。因此较重的第 $j$ 球产生第 $j$ 列向量，较轻时产生其相反向量；观察三次结果后即可唯一译码，所列 $24$ 个结果没有冲突。
若要处理十三球，必须用完 $3^3-1=26$ 个非零结果，每列与其相反数构成一对，十三列恰各代表一对。固定任意一行，在全部 $26$ 个非零三元向量中，该坐标非零的有 $2\cdot3^2=18$ 个，所以十三个代表列中这一行恰有九个非零项。九个球不可能均分到两盘，违背每次两盘球数相等。因此十三球不可能，十二球方案已达到这一类预定称量的上限。
\insight 信息计数只是必要条件；编码还必须满足实际称量的平衡约束。把轻重对应到互为相反的结果向量，可以把两者同时检查。


\par\Needspace{8\baselineskip}
\subsection*{小节练习}
\exercise{OA09}{受限取石游戏}
桌上有 $n\ge2$ 颗石子，两人轮流取，取走最后一颗者获胜。先手第一次可取 $1$ 至 $n-1$ 颗；以后每次至少取一颗，且不能超过对方上一次取走的数量，剩余不足时可以全部取走。双方均采取最优策略。哪些 $n$ 使先手必胜？给出必胜时的具体策略。

\exercise{OA13}{互素选数游戏}
两人从整数 $2,3,\ldots,30$ 中轮流选择一个尚未选过的数。每次新选的数必须与此前所有已选数分别互素；无数可选者输。双方均采取最优策略，谁能保证获胜？给出策略并证明。


\section{练习题}
\noindent\textbf{要点与方法.}\enspace
组合练习先辨认自由度：图中的未知边由连通结构约束，仿射排列由模意义下的可逆性约束，关联设计可按人或物品分别计数。随机化变换则先在单个原子条件下求分布，再按原子概率混合。

博弈练习要重新检查改变后的规则。取完者输会改变终止状态；收费会改变偏离收益，但不一定改变实际服务成本。求“全部”最优策略时，要给出概率或参数的完整范围，包括端点，并分别检验合作与竞争目标。
\par\smallskip



\exercise{AG2-1}{一个环的校验解}
图有顶点 $1,2,3,4$、四条环边 $12,23,34,41$ 和对角边 $13$。对角边值已知为 $1$，四个顶点的异或校验依次为 $(1,0,1,0)$。列出全部未知环边解。若四条未知边原来独立公平，满足校验的概率是多少？若再得知边 $12=0$，能否唯一恢复？

\exercise{AG2-2}{仿射皇后排列}
对 $5\times5$ 环面棋盘，在每行 $i$ 选择列 $j=ai+b\pmod5$，其中 $a,b\in\{0,1,2,3,4\}$。求满足 A14 全部条件的所有 $(a,b)$。再说明 $9\times9$ 时是否可能用任何方式安排，不限于仿射形式。

\exercise{AG2-3}{一个原子的均匀展开}
在 A28 的构造中，某原子 $x_0$ 满足 $F(x_0-)=3/8$、$\Prob(X=x_0)=1/4$。求 $U\mid X=x_0$ 的分布及 $\Prob(X=x_0,U>1/2)$。再证明无论如何选择确定性函数 $h$，$h(X)$ 都不可能服从连续均匀分布。

\exercise{AG2-4}{删除一名成员后的设计}
从 A17 的七人安排中删除一人，但保留全部物品，仍由其余原持有人持有。求恰由两人和恰由三人持有的物品数，并验证剩余六人仍各有三种、任意两人仍恰共有一种。用关联次数另列方程核对。

\pagebreak
\exercise{AG3-1}{环形搜查的全部均衡}
四个藏点按环排列为 $1,2,3,4$。甲选一个藏点，乙同时搜查一对相邻点，只能选 $12,23,34,41$；查到目标收益为 $1$，否则为零，双方为零和关系。求博弈值以及双方全部最优混合策略。

\exercise{AG3-2}{最后取完者输的取石游戏}
保留 A09 的取数集合 $\{1,3,4\}$，但改为拿走最后一颗者立即输。对正的剩余石子数，给出全部必败余数并证明。比较剩余 $20$ 颗时，本规则与 A09 规则下的必胜首步。

\exercise{AG3-3}{费用对拥挤均衡的调整}
A10 中取 $n=8$，向每位快服务器用户额外收取费用 $\delta\ge0$，慢服务器不收费。费用计入个人成本，但比较总服务成本时不计入转移支付。求使每个纯策略均衡都恰有三人使用慢服务器的全部费用值，并说明两个边界值的问题。

\exercise{AG3-4}{合作选址与竞争选址}
顾客在 $[0,1]$ 上均匀分布。两家服务商合作选择位置 $0\le a\le b\le1$，目标为最小化顾客到较近服务商的平均距离。求唯一的有序最优位置及最小距离。再与两家各自最大化需求份额、重合时平分需求的同时博弈比较。
